% GENERATED FILE. Edit canonical lessons, then run npm run generate:books.
% canonical-source-hash: fnv1a64:18dd1ed4b01765f5
% canonical-lessons: RU-C97-broom, RU-C97-washing-machine, RU-C97-sofa, RU-C97-wardrobe, RU-C97-saucepan

\chapter{Broom, Washing machine, Sofa, Wardrobe, Saucepan}
\label{ch:ru-broom-washing-machine-sofa-wardrobe-saucepan}
\hlchaptermodality{97}

\begin{chapteropening}
\textbf{By the end of this chapter:} \emph{I can say a broom, a washing machine, a sofa, a wardrobe and a saucepan.}

Complete the last lesson of chapter 97: I can say a broom, a washing machine, a sofa, a wardrobe and a saucepan.
\end{chapteropening}

\section[\texorpdfstring{\ru{веник} (vénik)}{веник (vénik)}]{\ru{веник} (vénik) — a broom}

\label{lesson:RU-C97-broom}



\begin{quote}
Before the new one: say the Russian for a comb, then the Russian for rubbish.
\end{quote}

\subsection*{You'll want to know: \ru{веник}}
\textbf{\ru{веник}} (\emph{vénik}) — \textquotedblleft{}a broom\textquotedblright{}.

\begin{cousinweb}
\textbf{\ru{веник}} — Old Russian \emph{venikŭ}, Czech \emph{věník}, \textquotedblleft{}a bundle\textquotedblright{}. Vasmer links it with \emph{\ru{венок}}, \textquotedblleft{}a wreath\textquotedblright{}, and \emph{\ru{вить}}, \textquotedblleft{}to twist\textquotedblright{}: a broom is a bound bundle of twigs.
\end{cousinweb}

\begin{grammarlens}[title={What you've built}]
One more word for this chapter.
\end{grammarlens}

\subsection*{Your turn}
\begin{itemize}
\raggedright
  \item \emph{Say it:} \emph{vénik}
  \item \emph{Say it:} \emph{vénik}, once more
  \item \emph{Say it:} say \emph{vénik}
  \item \emph{Recall:} say the Russian for a comb, then the Russian for rubbish, then say \emph{vénik} again
\end{itemize}

\begin{tcolorbox}[breakable,colback=teal!4,colframe=teal!35!black,title={Before you move on}]
What does \ru{веник} mean? (\textquotedblleft{}A broom\textquotedblright{}.) Say it once more.
\end{tcolorbox}
\section[\texorpdfstring{\ru{стиральная} \ru{машина} (stirálnaya mashína)}{стиральная машина (stirálnaya mashína)}]{\ru{стиральная} \ru{машина} (stirálnaya mashína) — a washing machine}

\label{lesson:RU-C97-washing-machine}



\begin{quote}
Before the new one: say the Russian for rubbish, then the Russian for a broom.
\end{quote}

\subsection*{You'll want to know: \ru{стиральная} \ru{машина}}
\textbf{\ru{стиральная} \ru{машина}} (\emph{stirálnaya mashína}) — \textquotedblleft{}a washing machine\textquotedblright{}.

\begin{cousinweb}
\textbf{\ru{стиральная} \ru{машина}} is \textquotedblleft{}washing machine\textquotedblright{}: \emph{\ru{стиральная}}, from \emph{\ru{стирать}}, \textquotedblleft{}to wash\textquotedblright{}, from \emph{\ru{тереть}}, \textquotedblleft{}to rub\textquotedblright{} (Latin \emph{terō}, \textquotedblleft{}I rub\textquotedblright{}, is a cousin), with \emph{\ru{машина}}, \textquotedblleft{}machine\textquotedblright{}, from German \emph{Maschine} (Vasmer).
\end{cousinweb}

\begin{grammarlens}[title={What you've built}]
One more word for this chapter.
\end{grammarlens}

\subsection*{Your turn}
\begin{itemize}
\raggedright
  \item \emph{Say it:} \emph{stirálnaya mashína}
  \item \emph{Say it:} \emph{stirálnaya mashína}, once more
  \item \emph{Say it:} say \emph{stirálnaya mashína}
  \item \emph{Recall:} say the Russian for rubbish, then the Russian for a broom, then say \emph{stirálnaya mashína} again
\end{itemize}

\begin{tcolorbox}[breakable,colback=teal!4,colframe=teal!35!black,title={Before you move on}]
What does \ru{стиральная} \ru{машина} mean? (\textquotedblleft{}A washing machine\textquotedblright{}.) Say it once more.
\end{tcolorbox}
\section[\texorpdfstring{\ru{диван} (diván)}{диван (diván)}]{\ru{диван} (diván) — a sofa}

\label{lesson:RU-C97-sofa}



\begin{quote}
Before the new one: say the Russian for a broom, then the Russian for a washing machine.
\end{quote}

\subsection*{You'll want to know: \ru{диван}}
\textbf{\ru{диван}} (\emph{diván}) — \textquotedblleft{}a sofa\textquotedblright{}.

\begin{cousinweb}
\textbf{\ru{диван}} — Vasmer thinks it more likely came from the West, French \emph{divan}, \textquotedblleft{}a couch\textquotedblright{}, than straight from its Turkish and Persian source \emph{divan}, \textquotedblleft{}a council, a council chamber\textquotedblright{} with its cushioned benches.
\end{cousinweb}

\begin{grammarlens}[title={What you've built}]
One more word for this chapter.
\end{grammarlens}

\subsection*{Your turn}
\begin{itemize}
\raggedright
  \item \emph{Say it:} \emph{diván}
  \item \emph{Say it:} \emph{diván}, once more
  \item \emph{Say it:} say \emph{diván}
  \item \emph{Recall:} say the Russian for a broom, then the Russian for a washing machine, then say \emph{diván} again
\end{itemize}

\begin{tcolorbox}[breakable,colback=teal!4,colframe=teal!35!black,title={Before you move on}]
What does \ru{диван} mean? (\textquotedblleft{}A sofa\textquotedblright{}.) Say it once more.
\end{tcolorbox}
\section[\texorpdfstring{\ru{шкаф} (shkaf)}{шкаф (shkaf)}]{\ru{шкаф} (shkaf) — a wardrobe, a cupboard}

\label{lesson:RU-C97-wardrobe}



\begin{quote}
Before the new one: say the Russian for a washing machine, then the Russian for a sofa.
\end{quote}

\subsection*{You'll want to know: \ru{шкаф}}
\textbf{\ru{шкаф}} (\emph{shkaf}) — \textquotedblleft{}a wardrobe, a cupboard\textquotedblright{}.

\begin{cousinweb}
\textbf{\ru{шкаф}} — most likely reshaped from older \emph{\ru{шкап}} under the influence of German forms such as \emph{Schaff}, Middle High German \emph{schaf}, \textquotedblleft{}a cupboard\textquotedblright{} (Vasmer).
\end{cousinweb}

\begin{grammarlens}[title={What you've built}]
One more word for this chapter.
\end{grammarlens}

\subsection*{Your turn}
\begin{itemize}
\raggedright
  \item \emph{Say it:} \emph{shkaf}
  \item \emph{Say it:} \emph{shkaf}, once more
  \item \emph{Say it:} say \emph{shkaf}
  \item \emph{Recall:} say the Russian for a washing machine, then the Russian for a sofa, then say \emph{shkaf} again
\end{itemize}

\begin{tcolorbox}[breakable,colback=teal!4,colframe=teal!35!black,title={Before you move on}]
What does \ru{шкаф} mean? (\textquotedblleft{}A wardrobe, a cupboard\textquotedblright{}.) Say it once more.
\end{tcolorbox}
\section[\texorpdfstring{\ru{кастрюля} (kastryúlya)}{кастрюля (kastryúlya)}]{\ru{кастрюля} (kastryúlya) — a saucepan}

\label{lesson:RU-C97-saucepan}



\begin{quote}
Before the new one: say the Russian for a sofa, then the Russian for a wardrobe.
\end{quote}

\subsection*{You'll want to know: \ru{кастрюля}}
\textbf{\ru{кастрюля}} (\emph{kastryúlya}) — \textquotedblleft{}a saucepan\textquotedblright{}.

\begin{cousinweb}
\textbf{\ru{кастрюля}} — first recorded in 1720; it probably came through German \emph{Kastrol}, \emph{Kasserolle}, from French \emph{casserole}, \textquotedblleft{}a saucepan\textquotedblright{} (Vasmer).
\end{cousinweb}

\begin{grammarlens}[title={What you've built}]
One more word for this chapter.
\end{grammarlens}

\subsection*{Your turn}
\begin{itemize}
\raggedright
  \item \emph{Say it:} \emph{kastryúlya}
  \item \emph{Say it:} \emph{kastryúlya}, once more
  \item \emph{Say it:} say \emph{kastryúlya}
  \item \emph{Recall:} say the Russian for a sofa, then the Russian for a wardrobe, then say \emph{kastryúlya} again
\end{itemize}

\begin{tcolorbox}[breakable,colback=teal!4,colframe=teal!35!black,title={Before you move on}]
What does \ru{кастрюля} mean? (\textquotedblleft{}A saucepan\textquotedblright{}.) Say it once more.
\end{tcolorbox}
